
\section{衍射现象，惠更斯-菲涅尔原理}

\begin{conclusion}[衍射现象与惠更斯-菲涅尔原理]
当光在传播中遇到障碍物或孔缝时，除了沿直线前进，还会绕过障碍物边缘进入几何阴影区，这种偏离直线传播的现象叫\textbf{衍射}。
常见的衍射图样有两类：圆孔衍射常呈\textbf{同心圆环}，单缝衍射常呈\textbf{明暗相间的条纹}。它们本质上都是波前被限制后，各部分子波相干叠加的结果。
\begin{itemize}
  \item 光的衍射图样沿着受限方向扩展。
  \item 光受限越厉害，衍射图样越展开，衍射效应越强。
\end{itemize}
惠更斯-菲涅尔原理指出：波前上每个面元 \(dS\) 都可看成一个次波源，观察点 \(P\) 处的振动，是波前 \(S\) 上各子波在该点产生振动的相干叠加，因此
\[
E(P)=\int_S \frac{C K(\theta)}{r}\cos\!\left(\omega t-\frac{2\pi r}{\lambda}\right)\mathrm dS .
\]
其中，\(S\) 表示 \(t\) 时刻的波前，\(dS\) 表示波阵面上面元（子波波源），\(r\) 为该面元到 \(P\) 点的距离，\(K(\theta)\) 为倾斜因子。按光源、孔缝与屏的相对距离，衍射可分为两类：\textbf{菲涅耳衍射}是光源、屏与缝相距有限远；\textbf{夫琅禾费衍射}是光源、屏与缝相距无限远。实际实验中常借助透镜把平行光照明和焦平面观察结合起来，从而实现夫琅禾费衍射。
\end{conclusion}

\begin{figure}[htbp]
\centering
\includegraphics[width=0.6\textwidth]{pdf-exercises/7.png}
\caption{菲涅尔衍射和夫琅禾费衍射的示意图}
\end{figure}
